\documentclass[11pt]{ctexart}
\usepackage[a4paper,margin=1in]{geometry}
\usepackage{graphicx}
\usepackage{xcolor}
\usepackage{hyperref}
\definecolor{refgray}{gray}{0.45}
\title{\textbf{市场最新观点汇总}\\[0.4em]{\large 全球市场风格切换，中国通胀传导断裂，AI与出口结构性重塑}}
\author{KC桌面 自动生成}
\date{260513}
\begin{document}
\maketitle
\section*{一页摘要}
\begin{itemize}
  \item 美国市场正式进入风险偏好模式，成长与动量因子主导，价值与质量因子退场，盈利超预期质量而非幅度成为定价核心。
  \item 中国4月PPI跳升至2.8\%，但上涨高度集中于上游能源与石化链条，下游消费品PPI仍为负，成本传导出现结构性断裂，中游制造业利润承压。
  \item 中国出口超预期增长14.1\%，AI相关产品（半导体、计算机设备）贡献近60\%增量，出口结构从成本优势转向技术供应链不可替代性。
  \item 央行一季度报告释放信号：降息窗口后移，政策重心从总量宽松转向结构性疏通，汇率管理可能进入“管理式升值”新阶段。
  \item AI正在重写网络攻击成本曲线，短期更有利于攻击者，网络安全支出正从合规成本转变为核心竞争力基础要素。
\end{itemize}
\section{宏观与利率：中国政策重心转向疏通，全球通胀传导分化}
\textbf{中国央行降息预期后移，政策重心从总量宽松转向结构性疏通；全球通胀呈现上游热、下游冷的结构性断层，利润分配向上游集中。}

\begin{itemize}
  \item 央行一季度报告明确短期降息概率极低，PPI转正与输入性通胀风险压缩降息空间，政策重点从“继续放水”转向“疏通水流”。
  \item 央行首次明确“引导隔夜市场利率围绕政策利率运行”，流动性管理从数量转向价格，银行间利率将更贴近政策利率。
  \item 4月PPI同比跳升至2.8\%，但上涨高度集中于石油与石化链条（贡献环比涨幅90\%），下游消费品PPI连续两月为负，成本传导断裂。
  \item 核心CPI和服务CPI走弱，房租跌幅创2023年以来最大，需求端并未真正回暖，通胀为供给侧冲击而非需求复苏。
  \item 央行报告重新出现“为实体经济营造稳定的汇率环境”表述，结合出口企业利润压力，人民币或进入“管理式升值”新阶段。
\end{itemize}
\begin{figure}[htbp]
\centering
\includegraphics[width=0.88\linewidth]{figures/fig_005.jpg}
\caption{Figure 1 - For interbank rate, the PBoC further specified to lead the overnight rate around the policy rate, compared with broad mentioning of short-term money market rates. This MPR didn't m}
\end{figure}
\begin{figure}[htbp]
\centering
\includegraphics[width=0.88\linewidth]{figures/fig_019.jpg}
\caption{Exhibit 1 - Exhibit 1 : Summary Table <table><tr><td></td><td>Apr-26</td><td>Mar-26</td><td>Feb-26</td><td>1Q26</td><td>4Q25</td></tr><tr><td>CPI, YoY \%</td><td>1.2</td><td>1.0</td><td>1.3</td}
\end{figure}
{\small\color{refgray}\textbf{References:} [R003] 央行一季度报告真正值得关注的信号：政策重心已从“降息”转向“疏通”，但市场可能低估了汇率管理的新含义; [R004] 市场真正低估的不是PPI跳升，而是通胀传导正在重塑产业链利润分配; [R007] 央行一季度报告最值得读的信号：宽松未变，但政策传导的“阀门”正在从数量转向价格; [R010] 市场真正低估的不是通胀本身，而是成本传导的结构性断裂; [R011] 中国通胀的真正信号不是价格涨了，而是涨价的逻辑变了; [R014] 市场真正低估的不是通胀，而是成本传导的结构性断层; [R020] 市场真正低估的不是通胀，而是上游涨价无法传导至下游的结构性断层}

\section{AI与科技：竞争焦点从模型能力转向商业化路径，半导体业绩分化加剧}
\textbf{中国大模型行业竞争进入下半场，商业化路径分化比模型能力差距更关键；半导体板块AI需求驱动盈利分化，估值逻辑从周期修复转向结构性重塑。}

\begin{itemize}
  \item 中国与美国在智能体能力上差距8-12个月，主因是数据质量和人才储备而非芯片短缺，视频生成能力已基本持平。
  \item 国产芯片目前仅覆盖不到40\%推理需求，但到2026年底有望提升至40-45\%，若7nm技术突破可升至50\%。
  \item DeepSeek降价是获取市场份额的战术动作，长期token价格将持续走低，商业化转向基于价值交付的服务型定价。
  \item 半导体板块1Q26营收增速跃升至30\%，但盈利分化远超营收：AI相关芯片（如澜起科技）净利润飙升61\%，而手机链和功率芯片承压。
  \item AI需求正在重塑估值逻辑：过去关注国产替代份额，未来关键变量是企业能否在AI驱动需求中占据不可替代生态位。
\end{itemize}
\begin{figure}[htbp]
\centering
\includegraphics[width=0.88\linewidth]{figures/fig_029.jpg}
\caption{Exhibit 8 - BofA GLOBAL RESEARCH \# Memory upcycles to continue According to BofA Global Memory Tech Team (report on 30 Apr), the memory “supercycle” is expected to continue into 2027-2028. Key}
\end{figure}
\begin{figure}[htbp]
\centering
\includegraphics[width=0.88\linewidth]{figures/fig_030.jpg}
\caption{Exhibit 9 - | 4Q27E | 11.3 | </details> Source: Company, BofA Global Research estimates BofA GLOBAL RESEARCH Exhibit 9: DRAM spot price softened in 2H March after an exceptionally strong rally}
\end{figure}
{\small\color{refgray}\textbf{References:} [R022] 半导体业绩分化的真正分水岭：AI需求正在重塑估值逻辑，而非只是周期修复; [R026] 中国大模型竞争进入下半场：商业化路径的分化比模型能力差距更值得关注}

\section{能源与大宗：上游涨价主导PPI，中下游利润承压}
\textbf{4月PPI跳升由能源和石化链条绝对主导，上游涨价无法有效传导至下游，中游制造业利润空间被系统性压缩。}

\begin{itemize}
  \item 石油和天然气开采PPI同比达29\%，煤炭采选PPI自2024年中以来首次转正，上游采掘和原材料价格暴涨。
  \item 制造业PPI仅微升至1.5\%，消费品PPI仍在-1.0\%负区间，汽车价格同比跌幅扩大至-1.2\%，下游缺乏议价能力。
  \item 能源价格传导虽存在政府管控，但汽油同比涨19.3\%已明显推高CPI，能源对CPI贡献约0.5个百分点。
  \item AI需求和全球铜价上涨推动有色金属价格走高，但反内卷政策下锂电池价格仅温和上涨，新能源车价格仍在下跌。
  \item 某投行上调2026年PPI预测至3.2\%，预计下半年触及5\%，但强调利润向上游集中，中下游毛利率将持续承压。
\end{itemize}
\begin{figure}[htbp]
\centering
\includegraphics[width=0.88\linewidth]{figures/fig_008.jpg}
\caption{Figure 2 - | Date | 2002 Dec | 2009 Dec | 2016 Sep | 2021 Jan | 2026 Mar | Consensus Fcst. | |------------|----------|----------|----------|----------|----------|-----------------| | Value | }
\end{figure}
\begin{figure}[htbp]
\centering
\includegraphics[width=0.88\linewidth]{figures/fig_026.jpg}
\caption{FIGURE 2 - | Apr-25 | -3.0 | -4.0 | -1.5 | | Oct-25 | -2.0 | -2.5 | -1.0 | | Apr-26 | 3.0 | 4.0 | -1.0 | </details> Source: Wind, BARC FIGURE 2. ...led by oil and gas extraction <details> <su}
\end{figure}
{\small\color{refgray}\textbf{References:} [R004] 市场真正低估的不是PPI跳升，而是通胀传导正在重塑产业链利润分配; [R010] 市场真正低估的不是通胀本身，而是成本传导的结构性断裂; [R011] 中国通胀的真正信号不是价格涨了，而是涨价的逻辑变了; [R014] 市场真正低估的不是通胀，而是成本传导的结构性断层; [R020] 市场真正低估的不是通胀，而是上游涨价无法传导至下游的结构性断层}

\section{出口与贸易：AI供应链重塑中国出口结构，进口激增揭示结构拐点}
\textbf{4月出口超预期增长14.1\%，AI相关产品贡献近60\%增量，出口结构从成本优势转向技术供应链不可替代性；进口激增揭示国内补库与大宗商品重估。}

\begin{itemize}
  \item 半导体出口同比近乎翻倍，自动数据处理设备增长47\%，两项合计贡献出口增速约8个百分点，占当月增量近60\%。
  \item 中国在全球关键AI相关产品出口中占比超30\%，全球AI资本开支周期启动时，中国是最大硬件制造枢纽。
  \item 进口同比增25.3\%，前四个月进口增速平均超出出口8个百分点，驱动力来自出口生产原材料需求、大宗商品涨价和国内补库。
  \item 出口目的地全线回暖：对美出口反弹11.3\%（低基数），对欧盟增13.4\%，对东盟增15.2\%，全球制造业PMI创四年新高。
  \item 出口结构升级正在改变就业传导效率，资本密集型行业每单位出口创造的就业远低于传统部门，内需疲弱格局短期难改。
\end{itemize}
\begin{figure}[htbp]
\centering
\includegraphics[width=0.88\linewidth]{figures/fig_009.jpg}
\caption{Figure 1 - Yi Xiong, Ph.D. Chief Economist +852-2203-6139 Figure 1: China' trade momentum continues in April <details> <summary>line</summary> | Year | China exports | China imports | |------}
\end{figure}
\begin{figure}[htbp]
\centering
\includegraphics[width=0.88\linewidth]{figures/fig_027.jpg}
\caption{FIGURE 1 - \# Details of April export data \textbackslash{}- By destination: Export growth improved broadly across major markets in April. Shipments to the US rebounded by 11.3\% y/y, partly reflecting a low }
\end{figure}
{\small\color{refgray}\textbf{References:} [R005] 市场真正低估的不是出口韧性，而是进口激增背后的结构拐点; [R012] 市场真正低估的不是出口韧性，而是出口拉动增长的结构性天花板; [R021] 出口的超预期不是周期反弹，而是中国供给结构被全球重新定价}

\section{地产与消费：一线城市现企稳信号，但复苏扩散仍遥远}
\textbf{中国房地产市场核心矛盾从“需求能否起来”转向“优质供给稀缺性重新定价”，一线城市二手房价格微幅回升，但复苏未扩散至更多城市。}

\begin{itemize}
  \item 4月85城二手房交易价格环比降0.4\%，跌幅较3月收窄，一线城市（上海、北京）二手房销售持续超预期，价格温和回升。
  \item AI供应链扩张带来的工业企业利润回升（1Q26同比增15\%）成为一线城市房价企稳新变量，计算机通信电子贡献8.2个百分点。
  \item 香港一手市场周末成交创2.5年新高，三个新盘全部售罄且定价高于二手14\%-31\%，市场进入“稀缺溢价”阶段。
  \item 餐饮行业4月回暖由清明假期脉冲驱动，五一假期数据明显降温，品牌分化加剧，补贴退坡后头部品牌面临压力测试。
  \item 宠物食品行业4月GMV同比增18\%，但成本端恶化（淀粉涨7\%、PET涨47\%），市场对品牌方盈利能力的定价可能过于乐观。
\end{itemize}
\begin{figure}[htbp]
\centering
\includegraphics[width=0.88\linewidth]{figures/fig_020.jpg}
\caption{Figure 1 - 7. Credit recommendations .... 18 8. Equity valuation summary 19 \# 1. Mainland China – Leading indicators Figure 1: Centraline secondary asking price index vs. NBS secondary home p}
\end{figure}
\begin{figure}[htbp]
\centering
\includegraphics[width=0.88\linewidth]{figures/fig_039.jpg}
\caption{Figure 2 - \# WHAT IS DIFFERENT THIS TIME? \# Breaking the negative feedback loop Over the past five years, there have been false alarms on property recovery, due to a weak job market, enterpri}
\end{figure}
{\small\color{refgray}\textbf{References:} [R016] 市场真正低估的不是需求，而是供给侧的再定价; [R023] 中国酒店业正在进入的不是需求复苏，而是供给侧的结构性定价权重构; [R024] 中国宠物食品市场真正的拐点，不是需求回暖，而是成本与竞争结构的双重挤压; [R025] 餐饮复苏的“弹簧效应”：四月回暖并非趋势反转，五一数据揭示了真正的分化; [R029] 中国房地产的财富效应正在被股市对冲，但奢侈品投资者不应过早乐观; [R031] 中国房地产的真正拐点，不是政策，是AI供应链带来的收入重构}

\section{区域市场：南非结构性改革催生赢家分化，港股流动性扩散创造alpha}
\textbf{南非市场机会从押注宏观复苏转向精选结构性赢家，金融板块确定性最高；港股流动性向中小盘扩散，南向资金结构变化创造选股机会。}

\begin{itemize}
  \item 南非改革议程整体在轨，物流领域是最大亮点（德班港吞吐量有望提升），但地方政府和水利改革滞后成为隐性风险。
  \item 金融板块是当前最高确信度主题，银行管理层对结构性改革带来的多年增长充满信心，同时严守回报率和信贷纪律。
  \item 港股活跃股票数量从2024年约230只激增至2025年约500只，流动性向更广泛标的扩散，中小市值股票可交易性显著提升。
  \item 南向资金中被动ETF持仓占比突破10\%，导致整体选股能力下降，但剔除ETF后主动型南向资金仍能产生稳定超额收益。
  \item 结合南向主动资金、全球最佳表现参与者及机构拥挤度信号的复合策略，年化超额回报达12\%，最大回撤仅6\%。
\end{itemize}
\begin{figure}[htbp]
\centering
\includegraphics[width=0.88\linewidth]{figures/fig_023.jpg}
\caption{Figure 1 - - Southbound investors from onshore China have been generally overweight in high-volatility names in the sectors such as information technology, whilst global participants tend to }
\end{figure}
\begin{figure}[htbp]
\centering
\includegraphics[width=0.88\linewidth]{figures/fig_024.jpg}
\caption{Figure 3 - A tale of two universes • After more than one decade of rapid development, Southbound has grown its coverage to most of the Hong Kong equity market's investable universe. - As of A}
\end{figure}
{\small\color{refgray}\textbf{References:} [R019] 港股真正的alpha不在总量，而在拆解南向资金的结构性变化; [R032] 南非市场的真正机会不在宏观反转，而在结构性改革催生的行业赢家分化}

\section{风险偏好：美股风格切换确认，网络安全成本曲线被AI重写}
\textbf{美国市场正式进入风险偏好模式，成长与动量因子主导，价值与质量因子退场；AI正在从根本上降低网络攻击成本，短期更有利于攻击者。}

\begin{itemize}
  \item 4月标普500涨超9\%，Big Tech涨超16\%，Momentum因子单月飙升18.6\%，1Q26 EPS惊喜率达18\%远超历史中位数5.2\%。
  \item 因子轮动信号明确：成长和动量上调至看多，价值和质量下调至中性/看空，市场押注风格切换而非盈利复苏。
  \item AI模型已展现出发现零日漏洞并自动生成利用代码的能力，攻击门槛断崖式下降，2025年美国网络攻击总损失约3000亿美元（≈1\%GDP）。
  \item 短期内AI更有利于攻击者而非防御者，因为攻击者只需找到一个入口，而防御者需守住所有出口，不对称性被急剧放大。
  \item 网络安全支出正从合规成本转变为核心竞争力基础要素，企业董事会和投资者需重新评估其战略性质。
\end{itemize}
\begin{figure}[htbp]
\centering
\includegraphics[width=0.88\linewidth]{figures/fig_001.jpg}
\caption{FIGURE 1 - - Risk assets posted an exceptional run in April, with the S\&P 500 up more than 9\% and Big Tech rallying by over 16\%, as Growth outperformed Value across both large and small caps.}
\end{figure}
\begin{figure}[htbp]
\centering
\includegraphics[width=0.88\linewidth]{figures/fig_013.jpg}
\caption{Exhibit 3 - A: Over one million cyber incidents were reported to the FBI's Internet Crime Complaint Center (IC3) in 2025. However, this probably understates the prevalence because cyberattacks}
\end{figure}
{\small\color{refgray}\textbf{References:} [R001] 市场真正定价的不是盈利复苏，而是风格切换的不可逆性; [R002] 市场真正在定价的不是风险偏好回归，而是对“估值叙事重构”的确认; [R008] 市场真正低估的不是攻击频率，而是AI正在重写网络攻击的成本曲线; [R009] S\&P 500 突破 8000 点并非妄想，但真正驱动下一轮上涨的并非盈利增长本身}

\section*{结语}
综合来看，当前市场主线清晰：全球风险偏好回归但结构分化加剧，中国通胀传导断裂与出口结构升级并存，AI正从技术概念演变为重塑产业利润分配和网络安全格局的核心变量。投资者需关注风格切换的不可逆性、成本传导的断层风险以及结构性改革催生的赢家分化。
\vfill
{\small\color{refgray}Personal reading notes and learning share only. Not investment advice.}
\end{document}
